\chapter{Metodologia}
\label{chap:methodology}

% Descrivere il sistema finale seguendo il percorso di dati e decisioni.
% Usare sorgenti e configurazioni delle run selezionate, senza raccontare
% la successione dei tentativi preliminari o anticipare i risultati.

\section{Formulazione del problema}
\label{sec:metodo-problema}

\subsection{Obiettivo di navigazione e costo energetico}
% Notazione, goal e compromesso fra successo, energia e durata;
% non attribuire al metodo garanzie di ottimalita.

\subsection{Informazione disponibile e osservabilit\`a parziale}
% Informazione locale e storia delle osservazioni; distinguere stato
% fisico completo e osservazione fornita alla policy.

\subsection{Assunzioni operative e vincoli}
% Formalizzare le assunzioni del capitolo 1, i limiti dei comandi
% e la dipendenza dal controllore locomotorio adottato.

\section{Piattaforma robotica e ambiente di simulazione}
\label{sec:metodo-piattaforma}

\subsection{Spot e controllore di locomozione}
% Robot e controllore preaddestrato che esegue i comandi di navigazione.

\subsection{Sensori e segnali disponibili}
% Camera, RGB-D, stato del corpo e segnali articolari realmente utilizzati.
% Distinguere segnali per il calcolo energetico e ingressi della policy.

\subsection{Simulazione mediante Isaac Sim e Isaac Lab}
% Framework, versione software e modello del robot impiegati.

\subsection{Modellazione dei terreni e delle interazioni fisiche}
% Geometria, materiali, contatti e proprieta fisiche. Disposizione delle
% mappe e condizioni specifiche di valutazione appartengono al capitolo 4.

\section{Architettura complessiva del sistema}
\label{sec:metodo-architettura}

\subsection{Flusso delle informazioni e moduli principali}
% Schema: RGB-D, feature, memoria/mappa, osservazione della policy,
% comandi di velocita, controllore locomotorio e feedback del robot.

\subsection{Collegamento fra policy di navigazione e controllore locomotorio}
% Separazione dei livelli di controllo; componenti addestrati e congelati.

\subsection{Frequenze operative e sincronizzazione}
% Frequenze di fisica, locomozione, percezione e navigazione;
% allineamento temporale di immagini, posa, azioni e misure energetiche.

\section{Stima del costo energetico}
\label{sec:metodo-energia}

\subsection{Potenza meccanica articolare}
% Somma dei valori assoluti dei prodotti fra coppia e velocita articolare;
% indicare la sorgente dei segnali nel simulatore.

\subsection{Integrazione temporale del lavoro meccanico}
% Metodo di integrazione e intervalli inclusi nella misura.

\subsection{Segnale energetico fornito alla policy e limiti della misura}
% Potenza filtrata nell'osservazione, costo nella reward ed energia cumulata;
% distinguere lavoro meccanico assoluto, lavoro netto e consumo elettrico.

\section{Rappresentazione percettiva e memoria dell'ambiente}
\label{sec:metodo-percezione}

\subsection{Acquisizione e preparazione dei dati RGB-D}
% Preparazione delle immagini, calibrazione geometrica e corrispondenza
% fra immagini, profondita e posa della camera.

\subsection{Estrazione delle feature visive mediante DINOv2}
% Modello, patch e feature; backbone congelato, senza classificazione
% esplicita dei materiali.

\subsection{Riduzione dimensionale mediante PCA}
% Trasformazione adottata e dimensione delle feature ridotte.

\subsection{Proiezione spaziale e descrittori geometrici}
% Riferimenti, proiezione dei dati e calcolo dei descrittori geometrici
% effettivamente usati, inclusi altezza, pendenza e rugosita.

\subsection{Costruzione della gridmap locale}
% Estensione, risoluzione, ordinamento e contenuto delle celle;
% distinguere regioni osservate e regioni prive di dati.

\subsection{Memoria delle osservazioni, validit\`a ed et\`a dei dati}
% Aggiornamento, interrogazione, trasformazioni spaziali, validita, eta
% e reset. Non attribuire alla memoria una mappa esplicita di consumi
% aggiornata online se il sistema non la implementa.

\section{Definizione del compito di reinforcement learning}
\label{sec:metodo-compito-rl}

\subsection{Spazio delle osservazioni}
% Tabella degli ingressi: significato, provenienza, dimensioni e unita.
% Collegare stato del robot, goal e rappresentazione locale.

\subsection{Spazio delle azioni e conversione in comandi}
% Azioni apprese e trasformazione nei comandi di velocita.

\subsection{Funzione di ricompensa}
% Definire i termini di progresso, energia, tempo, successo e fallimento
% e gli eventuali altri termini realmente presenti nella versione finale.

\subsection{Inizializzazione e condizioni di terminazione degli episodi}
% Reset, eventuale allineamento iniziale, successo, timeout, caduta
% e altre condizioni operative adottate.

\section{Architettura delle reti e fusione delle informazioni}
\label{sec:metodo-reti}

\subsection{Organizzazione e normalizzazione degli ingressi}
% Ordine delle variabili, trattamento dei dati mancanti e scaling.

\subsection{Rete della policy e fusione delle modalit\`a}
% Rami di elaborazione, MLP, attivazioni e distribuzione delle azioni
% secondo le architetture finali adottate.

\subsection{Critic e informazioni privilegiate}
% Reti/teste di valore e ingressi aggiuntivi del training, distinguendoli
% da quelli dell'actor impiegato durante l'esecuzione.

\section{Procedura di addestramento}
\label{sec:metodo-addestramento}

\subsection{Calibrazione dei componenti e preparazione dell'addestramento}
% Dati e procedura di calibrazione di PCA e normalizzazioni, componenti
% congelati e separazione fra dati di calibrazione e valutazione.

\subsection{Aggiornamento delle reti mediante PPO e SAC}
% Impiego concreto e particolarita implementative; rinviare ai principi
% del capitolo 2 senza ripeterne la trattazione.

\subsection{Generazione delle missioni e curriculum}
% Campionamento di start/goal, fasi e criteri di avanzamento.
% Distinguere gli strumenti di generazione delle missioni dagli input
% effettivamente forniti all'actor.

\subsection{Configurazioni di addestramento adottate}
% Configurazioni finali degli scenari: iperparametri, parallelismo, budget
% e criteri di selezione. Non ricostruire la cronologia delle prove e non
% implicare un confronto controllato PPO/SAC se non viene svolto.

\section{Esecuzione della policy e integrazione sul robot}
\label{sec:metodo-esecuzione}

\subsection{Pipeline di inferenza}
% Caricamento della policy e componenti necessari all'esecuzione.

\subsection{Gestione delle osservazioni e dei comandi durante l'esecuzione}
% Acquisizione dei segnali, aggiornamento della mappa, scaling e invio
% dei comandi, in coerenza con la pipeline di addestramento.

\subsection{Integrazione con la piattaforma reale}
% Da completare quando l'implementazione reale sara definita.
% Descrivere il sistema effettivamente utilizzato; i risultati sono nel
% capitolo 4. Non presentare il trasferimento come gia validato.
